\chapter*{Résumé}
\addcontentsline{toc}{chapter}{Résumé}

Ce projet, réalisé dans le cadre du module \emph{Internet des objets}, porte sur la conception d'un système
intelligent de supervision d'un véhicule capable de détecter et de prendre en compte les obstacles présents sur
son trajet. Le système associe trois briques complémentaires.

La première est un \textbf{prototype matériel} : une voiture miniature construite autour d'une carte Arduino UNO,
d'un pont en H L298N, de deux moteurs à courant continu, de deux capteurs infrarouges qui suivent une piste
tracée au ruban noir et d'un capteur ultrasonique HC-SR04. Une caméra ESP32-CAM autonome, alimentée par une
batterie externe, diffuse en Wi-Fi les images prises du point de vue de la voiture.

La deuxième brique est un ensemble de \textbf{modèles de vision par ordinateur} de la famille YOLO26
(Ultralytics). Pour l'environnement réel, un modèle préentraîné sur COCO reconnaît les usagers de la route
(piétons, cyclistes, véhicules, animaux) et un second modèle, affiné pendant huit heures sur 7\,894 images de
caméras embarquées, détecte les anomalies routières (nids-de-poule, cônes, barrières, panneaux de chantier). Pour
le prototype, un modèle YOLO26n a été entraîné sur les images de notre propre piste, annotées automatiquement par
suivi vidéo avec SAM~2, afin de reconnaître les deux obstacles du circuit.

La troisième brique est l'application web \textbf{SmartCar Ops}, développée avec Dash et Plotly et déployée avec
Docker. Elle exécute les modèles en temps réel (ONNX Runtime), estime la distance et la gravité de chaque
obstacle, lui attribue un niveau de priorité, le place sur une carte et en tient compte pour proposer
l'itinéraire le plus sûr. Huit pages (vue d'ensemble, suivi, caméra, planification, obstacles, statistiques,
modèles d'IA, périphériques et journaux) permettent de superviser deux environnements : un environnement
\emph{monde réel} (flotte simulée sur le réseau routier du Quartier Latin à Paris, avec de vraies vidéos de route
analysées par les modèles) et l'environnement \emph{prototype} de la voiture miniature.

Sur l'ensemble de test, le détecteur d'anomalies routières atteint un mAP@50 de 0,53 en temps réel (0,57 au
meilleur réglage), et le modèle du prototype détecte tous les obstacles des images de test (mAP@50 de 0,995),
résultat à nuancer par la petite taille de cet ensemble.

\vspace{0.5cm}
\noindent\textbf{Mots-clés :} Internet des objets, Arduino, ESP32-CAM, vision par ordinateur, détection d'objets,
YOLO26, Ultralytics, SAM~2, Dash, Plotly, Docker, planification d'itinéraire.
